\documentclass[a4paper]{article}


\usepackage[teacher]{Lecture-TD}

\lhead{Travaux dirigés - Performances des accumulateurs}
\lfoot{IUT d'Orsay - BUT 2A - 2024/2025} \cfoot{
\includegraphics[height=0.5cm]{cc-by-nc-sa.png}}  \rfoot{Jean-François Olivieri \quad \thepage}

\begin{document}

\begin{center}
\Large \textbf{TD n$^\circ ~ 2$} \\
Performances des accumulateurs
\end{center}

\objectives{
\begin{itemize}
\item Savoir quantifier la capacité théorique d'une batterie ;
\item Savoir interpréter des valeurs de potentiels rédox ou des valeurs de force électromotrice ;
\item Savoir prédire les performances théoriques d'un dispositif de stockage d'énergie ;
\item Savoir interpréter une courbe de charge/décharge ;
\item Savoir interpréter les paramètres pouvant influencer les courbes de charge/décharge : température, vitesse de charge ou de décharge, nombre de cycle ;
\item Savoir calculer et interpréter le SoC, SoD et SoH.
%\item Savoir calculer un rendement en charge et en puissance d'une batterie.
\end{itemize}
}

\cteacher{Dans ce TD, il s'agit avant tout de les aider à faire un point sur les notions élémentaires : charge, courant, tension, force électromotrice, travail électrique. \\
Dans la première partie, il s'agit de leur montrer qu'il est possible de calculer certaines grandeurs théoriques associées aux performances de l'accumulateur. \\
Pour la seconde partie, nous avons montré en CM l'inégalité suivante : 
\begin{align*}
\langle U_c \rangle > e > \langle U_d \rangle
\end{align*}
avec $\langle U_c \rangle$ (resp. $\langle U_d \rangle$) la tension moyenne en charge (resp. en décharge) et $e$ la force électromotrice. \\
Cependant, les paramètres cinétiques à l'origine de l'hystérèse n'ont pas été discutés. Vous pouvez leur faire un petit point sur les surtensions et la chute ohmique principalement, si nécessaire, dans vos explications.
Attention à les aider sur les grandeurs telles que le SoC, le DoD ou le SoH sur lesquels j'ai passé assez peu de temps en CM. \\
Les exercices I.1, I.2 et II.2 sont à traiter en priorité. Vous pouvez leur laisser le II.1 à faire seuls.
}


%\cteacher{
%Les sources de f.e.m. : galvanique, concentration et température peuvent être abordées en s'appuyant sur la relation de Nernst. Ne traiter que la pile de concentration. Une analogie, non traitée en CM, est à établir à la fin de l'exercice : batterie Li-ion vs pile de concentration. Bien entendu, il y a de grosses limites à cette analogie, mais il se trouve que l'on retrouvera, pour la f.e.m., des expressions mathématiques similaires.
%}

\section{Performances théoriques maximales}


\subsection{Comparaison des capacités spécifiques massiques de quelques électrodes courantes des batteries \ce{Li}}
%from Culture science chimie

\setcounter{equation}{0}
\setcounter{Q}{0}

La capacité spécifique massique est la charge électrique pouvant être délivrée par gramme de matière active, c'est-à-dire par gramme du matériau hôte. Dans le système international, on l'exprime en \si{\coulomb\per\gram} mais en pratique, on l'utilise plutôt en \si{\milli\ampere\hour\per\gram}.

\Q \label{Q:11-01} Définir un ampère-heure [\si{\ampere\hour}] et établir le facteur de conversion avec le coulomb [\si{\coulomb}].
\solution{Un ampère-heure est une unité qui représente la quantité de charge délivrée par un dispositif débitant 1 ampère pendant 1 heure. Or, \SI{1}{\hour} équivaut à \SI{3600}{\second}, on peut donc en déduire que : 
\begin{align*}
\SI{1}{\ampere\hour} = \SI{3600}{\coulomb}
\end{align*}}

\Q \label{Q:11-02} On va considérer une réaction rédox modèle de la forme : 
\begin{align*}
\ce{M + z Li^+ + z e^- = MLi_z(s)}
\end{align*}
où $M$ représente le matériau hôte du lithium.

\'{E}tablir que la capacité spécifique massique s'exprime en \si{\milli\ampere\hour\per\gram} et peut être estimée par la relation :  
\begin{align*}
C_\text{Hote} = \frac{z F}{3.6 M_{Hote}}
\end{align*}

où $z$ est le nombre d'électrons échangé, $F$ la constante de Faraday et $M_{Hote}$ la masse molaire du matériau hôte.

\solution{
On va considérer que l'on échange $n_e$ moles électrons et de lithium avec un matériau hôte de masse $m_{hote}$, c'est-à-dire que le matériau peut accueillir au maximum $n_e$ moles de lithium dans sa structure ! Dressons un tableau d'avancement de la situation : 
\begin{center}\begin{tabular}{cccccccc}
\si{\mole} & \ce{M(s)} & \ce{ + } & \ce{z Li^+} & \ce{ + } & \ce{z e^-} & \ce{ = } & \ce{MLi_z(s)} \\
$t=0$ & $\frac{m_{hote}}{M_{hote}}$ & & $n_e$ & & $n_e$ & & 0\\
$t=\infty$ & $\frac{m_{hote}}{M_{hote}} - \xi_{\infty}$ & & $n_e - z\xi_{\infty}$ & & $n_e - z \xi_{\infty}$ & & $\xi_\infty$\\
\end{tabular}\end{center}
Lorsque l'échange d'électrons est terminé, tous les électrons échangeables ont été transférés, donc : 
\begin{align*}
\begin{cases}
\frac{m_{hote}}{M_{hote}} - \xi_{\infty} &= 0 \\
n_e - z\xi_{\infty} = 0
\end{cases} \implies \frac{m_{hote}}{M_{hote}} = \frac{n_e}{z} \implies n_e = \frac{z m_{hote}}{M_{hote}}
\end{align*}
Ainsi, on a une relation simple entre la structure hôte et le nombre d'électrons échangé. Or, on sait exprimer la quantité d'électrons échangés en quantité de charges échangées $Q$ : 
\begin{align*}
Q = F n_e = \frac{z F m_{hote}}{M_{hote}} \quad [\si{\coulomb}]
\end{align*}
La capacité spécifique massique de cette structure correspond au nombre d'électrons échangés par gramme de matériau, que l'on peut exprimer en \si{\coulomb\per\gram} puis en \si{\ampere\hour\per\gram} grâce à la Q.\ref{Q:11-01} :
\begin{align*}
C_{hote} &= \frac{Q}{m_{hote}} = \frac{z F }{M_{hote}} \quad [\si{\coulomb\per\gram}] \\
&= \frac{z F}{3.6 M_{hote}} \quad [\si{\milli\ampere\hour\per\gram}]
\end{align*}
}

\Q Dans le cas des matériaux lithiés suivants, donner la demi-équation rédox de réduction du cation lithium en lithium métallique. En déduire la capacité spécifique et compléter le tableau :

\begin{center}
\begin{tabular}{ccccccc}
\toprule
Matériau & \ce{Li(s)} & \ce{LiC6(s)} & \ce{Li15Si4(s)} & \ce{LiCoO2(s)} & \ce{LiFePO4(s)} & \ce{LiNi_{1/3}Co_{1/3}Mn_{1/3}O2} \\
\midrule
Masse molaire [\si{\gram\per\mole}] & 6.94 & 12.0 & 28.1 & 97.9 & 158 & 96.5 \\
\midrule
Capacité [\si{\milli\ampere\hour\per\gram}] & & & & & & \\
%Capacité [\si{\milli\ampere\hour\per\gram}] & 3960 & 372 & 3600 & 274 & 170 & 279 \\
\bottomrule
\end{tabular}
\end{center}

Commenter le résultat. \emph{Aide : Vous penserez à choisir une st{\oe}chiométrie compatible avec l'équation modèle étudiée à la Q. \ref{Q:11-02}.}

\solution{On peut écrire toutes les équations rédox en réduction, suivant le modèle de l'énoncé : 
\begin{align*}
& \ce{Li^+ + e^- = Li} & z = 1\\
& \ce{\frac{1}{6} Li^+ + \frac{1}{6} e^- + C = \frac{1}{6} LiC6} & z = 1/6\\
& \ce{\frac{15}{4} Li + \frac{15}{4} e^- + Si = \frac{1}{4} Li_{15}Si4} & z = 15/4\\
& \ce{Li^+ + e^- + CoO2 = LiCoO2} & z = 1 \\
& \ce{Li^+ + e^- + FePO4 = LiFePO4} & z = 1 \\
& \ce{Li^+ + e^- + Ni_{1/3}Co_{1/3}Mn_{1/3}O2 = LiNi_{1/3}Co_{1/3}Mn_{1/3}O2} & z=1
\end{align*}
On peut alors compléter le tableau proposé : 
\begin{center}
\begin{tabular}{ccccccc}
\toprule
Matériau & \ce{Li(s)} & \ce{LiC6(s)} & \ce{Li15Si4(s)} & \ce{LiCoO2(s)} & \ce{LiFePO4(s)} & \ce{LiNi_{1/3}Co_{1/3}Mn_{1/3}O2(s)} \\
\midrule
Masse molaire [\si{\gram\per\mole}] & 6.94 & 12.0 & 28.1 & 97.9 & 158 & 96.5 \\
\midrule
Capacité [\si{\milli\ampere\hour\per\gram}] & 3860 & 372 & 3580 & 274 & 170 & 279 \\
\bottomrule
\end{tabular}
\end{center}
On observe qu'en général, les matériaux qui accueillent le lithium n'atteignent pas \SI{10}{\percent} des capacités du lithium métallique. Cependant, une exception intéressante apparaît dans la liste : \ce{Li15Si4}. En effet, les matériaux au silicium sont aujourd'hui des candidats prometteurs pour leur excellent rapport Li/Si (=très nombreux sites d'insertion) qui permettent d'approcher des performances en capacité du lithium métallique ! \\
\textbf{Nota Bene :} La capacité spécifique ne dépend pas de la "nature de la structure cristallographique" du matériau.}

\paragraph*{Données}
\begin{itemize}
\item Constante de Faraday : $F = 96485 ~ \si{\coulomb\per\mole}$ 
\end{itemize}

\begin{comment}
\begin{itemize}
\item Masses molaires des éléments : 
\begin{center}
\begin{tabular}{cccccccc}
Lithium & Carbone & Oxygène & Silicium & Cobalt & Fer & Manganèse & Phosphore \\
6.94 & 12.0 & 16.0 & 28.1 & 59.0 & 55.9 & 55.0 & 31.0
\end{tabular}
\end{center}
\end{itemize}
\end{comment}


\subsection{Energie théorique maximale}


\setcounter{equation}{0}
\setcounter{Q}{0}

A l'exercice précédent, nous avons vu qu'il était possible de calculer les performances théoriques maximales d'un matériau d'électrode.

En particulier, ces performances sont mesurées à l'aide d'un accumuluteur présentant un pôle $\ominus$ (souvent appelé anode dans la littérature) en lithium métallique : 

\begin{align*}
\oplus \quad Li_x M(s) | Li^+ || Li^+ | Li(s) \quad \ominus
\end{align*}

où $M$ désigne un matériau d'életrode donné.

Afin de comparer les performances des électrodes du pôle $\oplus$ (souvent appelé cathode dans la littérature), nous allons proposer les potentiels suivants : 

\begin{center}
\begin{tabular}{ccccccc}
\toprule
Matériau & \ce{Li(s)} & \ce{LiC6(s)} & \ce{Li15Si4(s)} & \ce{LiCoO2(s)} & \ce{LiFePO4(s)} & \ce{LiNi_{1/3}Co_{1/3}Mn_{1/3}O2(s)} \\
\midrule
$E_{\oplus}^{eq}$ [$\si{\volt\per{Li^+/Li}}$] & 0 & 0.3 & 0.6 & 4.1 & 3.3 & 3.9  \\
\bottomrule
\end{tabular}
\captionof{figure}{Potentiel électrochimique d'équilibre à \SI{298}{\kelvin} de différents matériaux du pôle $\oplus$}
\end{center}



\Q Qu'appelle-t-on une pseudo-référence ? Qui joue ce rôle ici ? Que signifie $\si{\volt\per{Li^+/Li}}$ ?

\solution{
Les électrodes de pseudo-référence sont des électrodes fabriquée à partir de matériaux dont le potentiel électrochimique à l'aquilibre est stable ou dont la variation est prévisible.\\
L'électrode au lithium métallique est considérée ici comme une pseudo-référence dont le potentiel est supposé constant et valoir approximativement : 
\begin{align*}
E_\oplus^{eq} \approx - 3.04 ~ \si{\volt\per{ESH}}
\end{align*}
L'unité $\si{\volt\per{Li^+/Li}}$ signifie que l'on considère le potentiel de l'électrode $\oplus$ par rapport au couple \ce{Li^+}/\ce{Li} à l'électrode $\ominus$ : 
\begin{align*}
E_\oplus^{eq} [\si{\volt\per{Li^+/Li}}] &= E_\oplus^{eq} [\si{\volt\per{ESH}}] - E_\ominus^{eq} [\si{\volt\per{ESH}}] \\
&=E_\oplus^{eq} [\si{\volt\per{ESH}}] - E^{eq} (\ce{Li^+} / \ce{Li})
\end{align*}
}

\Q En déduire la f.e.m. dans ces cellules Li-metal. 
\solution{Par conséquent, la force électromotrice dans ces cellules est approximativement égale au potentiel du pôle $\oplus$ :
\begin{center}
\begin{tabular}{ccccccc}
\toprule
Matériau & \ce{Li(s)} & \ce{LiC6(s)} & \ce{Li15Si4(s)} & \ce{LiCoO2(s)} & \ce{LiFePO4(s)} & \ce{LiNi_{1/3}Co_{1/3}Mn_{1/3}O2(s)} \\
\midrule
$f.e.m.$ [$\si{\volt}$] & 0 & 0.3 & 0.6 & 4.1 & 3.3 & 3.9  \\
\bottomrule
\end{tabular}
\end{center}
}

\Q On observe deux catégories de matériaux à partir des valeurs de f.e.m. Quelles sont-elles ?  Donner un sens énergétique à la f.e.m.
\solution{On peut classer les matériaux en fonction de leur proximité en potentiel avec l'électrode au lithium. Les électrodes $\ominus$ (anode en décharge) doivent avoir un potentiel le plus proche possible du lithium, tandis que les électrodes $\oplus$ (cathode en décharge) doivent avoir le potentiel le plus éloigné possible du lithium métallique : 
\begin{center}
\begin{tabular}{ccccccc}
\toprule
Matériau & \ce{Li(s)} & \ce{LiC6(s)} & \ce{Li15Si4(s)} & \ce{LiCoO2(s)} & \ce{LiFePO4(s)} & \ce{LiNi_{1/3}Co_{1/3}Mn_{1/3}O2(s)} \\
\midrule
$f.e.m.$ [$\si{\volt}$] & 0 & 0.3 & 0.6 & 4.1 & 3.3 & 3.9  \\
Rôle & $\ominus$ & $\ominus$ & $\ominus$ & $\oplus$ & $\oplus$ & $\oplus$ \\
\bottomrule
\end{tabular}
\end{center}
} 

\Q Après avoir exprimé le travail électrique, calculer l'énergie théorique maximale stockable dans ces cathodes.
\solution{
En convention récepteur (resp. générateur), la puissance électrique reçue (resp. cédée) vaut : 
\begin{align*}
P = U \times I
\end{align*}
Le travail électrique reçu vaut : 
\begin{align*}
W' = \int P(t) dt = \int U(Q) dQ
\end{align*}
Au maximum, en décharge, la tension vaut la force électromotrice. L'énergie théorique maximale s'écrit donc :
\begin{align*}
ETM [\si{\milli\watt\hour\per\gram}] = e [\si{\volt}] \times CS_m [\si{\milli\ampere\hour\per\gram}]
\end{align*}
On peut alors calculer l'énergie théorique maximale, siglée ETM, pour chacune des électrodes :
\begin{center}
\begin{tabular}{ccccccc}
\toprule
Matériau & \ce{Li(s)} & \ce{LiC6(s)} & \ce{Li15Si4(s)} & \ce{LiCoO2(s)} & \ce{LiFePO4(s)} & \ce{LiNi_{1/3}Co_{1/3}Mn_{1/3}O2(s)} \\
\midrule
$f.e.m.$ [$\si{\volt}$] & 0 & 0.3 & 0.6 & 4.1 & 3.3 & 3.9  \\
Capacité [\si{\milli\ampere\hour\per\gram}] & 3860 & 372 & 3580 & 274 & 170 & 279 \\
ETM [\si{\milli\watt\hour\per\gram}] & $0$ & $112$ & $2148$ & $1123$ & $561$ & $1088$ \\
\bottomrule
\end{tabular}
\end{center}
La valeur d'un matériau comme électrode au pôle $\oplus$ doit tenir compte de deux choses : de la densité d'énergie stockée, mais aussi de la qualité de cette énergie.
}

\cteacher{Attention, ce diagramme est analogue au diagramme de Ragone mais n'en est pas un. Dans le diagramme de Ragone, on représente l'énergie en fonction de la puissance. Vous pouvez leur préciser cette erreur dans l'énoncé, bien qu'elle ne soit pas critique.}
Le diagramme de Ragone permet de comparer graphiquement les performances de dispositifs de stockage d'énergie (condensateurs, accumulateurs, pile à combustible...) ou de matériaux au sein d'un même dispositif, comme c'est le cas ici pour l'accumulateur au lithium.

\begin{figure}[H]
\subfloat[\label{fig:01a}]{\includegraphics[height=5cm]{Ex37/Fig01.png}}\subfloat[\label{fig:01b}]{\includegraphics[height=5cm]{Ex37/Fig02.png}}
\caption{Potentiel moyen de fonctionnement de quelques matériaux d'électrode négative et positive en fonction des valeurs de capacité accessibles expérimentalement ($<$ aux valeurs théoriques). Du coté anodique $\ominus$, les matériaux d’intercalation (spinelle LTO : \ce{Li4Ti5O12}, graphite, carbone dur), les matériaux de conversion (oxydes comme \ce{SnO2} et fluorures, phosphures et sulfures), les alliages (étain, germanium, silicium) et le lithium métallique. Du coté cathodique $\oplus$, les matériaux lamellaires [LTS : \ce{LiTiS2}, LCO : \ce{LiCoO2} et LCO partiellement substitués : NMC (\ce{Li(Ni,Mn,Co)O2}) et NCA (\ce{Li(Ni,Co,Al)O2})], les spinelles (LMO : \ce{LiMn2O4}, et LMO partiellement substitué LNMO) et les matériaux polyanioniques (olivine - LFP : \ce{LiFePO4}, LCP : \ce{LiCoPO4} et tavorite LFSF : \ce{LiFeSO4F})}
\end{figure}%http://reperageterrestre.free.fr/coordonnees.html


\Q Où trouve-t-on les matériaux de haute et basse densité d'énergie dans les graphes précédents ?
\solution{
Les matériaux de haute (resp. basse) densité d'énergie se trouve en haut (resp. bas) à droite (resp. gauche) du graphique de Ragone.
}
\Q Une grandeur est manquante pour caractériser pleinement une cellule électrochimique. Laquelle ? Comment pourrait-on l'y ajouter ?
\solution{
Bien que les batteries ne soient pas des dispositifs de puissance, les batteries doivent être capable de débiter une certaine énergie par unité de temps. \\
Pour représenter ces performances en puissance, on utilise le diagramme de Ragone qui représente l'énergie en fonction de la puissance, le plus souvent en échelle log-log :
\begin{align*}
\log ETM = \log P_s + \log \tau
\end{align*}
avec $\tau$ le temps caractéristique de décharge de l'accumulateur.
\begin{center}
\begin{minipage}{0.48\textwidth}
\centering
\includegraphics[width=0.5\textwidth]{Ex37/Corr01}
\end{minipage}\hfill\begin{minipage}{0.48\textwidth}
\centering
\includegraphics[width=0.5\textwidth]{Ex37/Corr02}
\end{minipage}
\end{center}
}

\section{Performances expérimentales}

\subsection{Comparaison des performances de quelques technologies de batterie}

\setcounter{equation}{0}
\setcounter{Q}{0}

Les performances d'une batterie sont jugées en premier lieu à la densité massique d'énergie. En effet, on cherche en général à stocker un maximum d'énergie dans un dispositif léger qui prend le minimum de place.%avec un minimum d'encombrement.
\begin{center}
\begin{adjustbox}{width=0.95\textwidth}
\begin{tabular}{cccccc}
    \toprule
     \textbf{Technologie} & \textbf{Année de développement} & \textbf{\'{E}nergie massique} & \textbf{\'{E}nergie volumique} & \textbf{Tension} & \textbf{Cyclabilité}\\
     Unité & & \si{\watt\hour\per\kilo\gram} & \si{\watt\hour\per\liter} & \si{\volt}/au couple \ce{Li^+}/\ce{Li} & $\emptyset$ \\
     \midrule
     Plomb & 1859 & 35 & 70 & 2.0 & 200-250\\
     Nickel-Cadmium & 1900 & 60 & 100 & 1.2 & 400-500\\
     Nickel-Métal-Hydrure & 1975 & 80 & 240 & 1.2 & 400-500\\
     Lithium-Ion & 1978 & 250 & 400 & 3.6 & > 1000\\
     \bottomrule
\end{tabular}   
\end{adjustbox}
\captionof{table}{\small \'{E}volution de l'ordre de grandeur de l'énergie massique, volumique, de la tension ainsi que de la cyclabilité d'une batterie en fonction de la technologie.}
\label{tab:battery_time}
\end{center}

\Q Définir la densité massique d'énergie et la densité volumique d'énergie. Définir le wattheure et rappeler sa relation aux unités du système international. 
\solution{La densité massique/volumique d'énergie électrique délivrable par une batterie correspond au produit : 
\begin{align*}
\text{tension de la cellule} \times \text{capacité massique/volumique}
\end{align*}
Cette grandeur permet donc de condenser deux informations cruciales : le nombre d'électrons échangeables et le "potentiel" électrique de chacun de ces électrons ! \\
Le joule \si{\joule} est l'unité de l'énergie dans le système international. Le watt-heure représente l'énergie délivrée par un dispositif électrique fournissant une puissance électrique de \SI{1}{\watt} pendant \SI{1}{\hour}. Or, il y a \SI{3600}{\second} dans \SI{1}{\hour}, d'où  
\begin{align*}
\SI{1}{\watt\hour} = \SI{3600}{\joule}
\end{align*}}

\Q Déterminer la capacité spécifique massique et volumique moyenne de chacune des technologies.
\solution{On note $CS_m$ (resp. $CS_V$) les capacités spécifiques massiques (resp. volumique) de chacune de ces technologies.
\begin{center}
\begin{adjustbox}{width=0.95\textwidth}
\begin{tabular}{cccc}
    \toprule
     \textbf{Technologie} & \textbf{Année de développement} & \textbf{Capacité massique} & \textbf{Capacité volumique} \\
     Unité & & \si{\ampere\hour\per\kilo\gram} & \si{\ampere\hour\per\liter} \\
     \midrule
     Plomb & 1859 & 17.5 & 35 \\
     Nickel-Cadmium & 1900 & 50 & 84\\
     Nickel-Métal-Hydrure & 1975 & 67 & 200\\
     Lithium-Ion & 1978 & 70 & 111\\
     \bottomrule
\end{tabular}   
\end{adjustbox}
\end{center}}

\Q Commenter les valeurs de la Tab.\ref{tab:battery_time}. Quel est l'intérêt de la technologie Li-ion par rapport à d'autres technologies ?
\solution{On peut soulever les avantages de la batterie Li : 
\begin{itemize}
\item l'énergie massique la plus grande ;
\item l'énergie volumique la plus grande ;
\item la meilleure f.e.m. ;
\item la meilleure cyclabilité.
\end{itemize}
Les batteries lithium sont donc parfaitement adaptées aux dispositifs portatifs nécessitant de grandes quantités d'énergies (\emph{par ex.} les portables).  \\
Il est à noter toutefois que la technologie Ni-M-H présente des capacités spécifiques volumiques et massiques équivalentes, voire plus intéressantes, que la technologie lithium ce qui peut offrir des performances intéressantes pour les applications nécessitant de faibles capacités en énergie (\emph{par ex.} les terminaux de paiement bancaire, les technologies NFC, les accessoires d'Airsoft, certains capteurs). 
}

\Q Rappeler ce qu'est la qualité de l'énergie stockée et expliquer l'intérêt des batteries Li. Expliquez-en l'origine.

\solution{La qualité d'énergie est une façon d'exprimer l'origine de la puissance électrique qui est soit potentielle (venant de la f.e.m.), soit cinétique (venant du courant électrique). En effet, dans un modèle de Thévenin simple (cf cours), la puissance électrique est donnée par : 
\begin{align*}
P = e I - R I^2
\end{align*}
Un dispositif présentant une grande tension à vide $e$ présente l'avantage de délivrer moins de courant pour une même puissance électrique, ce qui a pour conséquence :
\begin{itemize}
\item de diminuer les déperditions d'énergie, sous forme thermique, dû aux phénomènes résistifs (terme en $-R I^2$) ;
\item d'augmenter la durée de décharge (c.-à.-d. que la batterie se décharge plus lentement).
\end{itemize}}

\Q Mettre en relation ces données avec la place de l'élément Li au sein du tableau périodique.
\solution{Le lithium se situe en haut à gauche du tableau périodique ($3^e$ élément du tableau). C'est un élément métallique, parmi les plus électropositifs (explique la f.e.m.), de faible poids moléculaire (explique les capacités massiques) et de faible nombre quantique principal (soit un élément de petite taille, ce qui explique les capacités volumiques).}

\subsection{Caractéristiques d'une pile commerciale}
%Panasonic

\setcounter{equation}{0}
\setcounter{Q}{0}


Panasonic Energy{\copyright} commercialise l'accumulateur ML2020, appelé \emph{Coin-type Manganese Rechargeable Lithium Batteries}, dont les spécifications techniques suivantes sont fournies : 
\begin{center}
\begin{minipage}{0.70\textwidth}
\begin{tabular}{cc}
\toprule
Spécifications & Valeurs prises \\
\midrule
Tension de charge & \SI{2.8}{\volt}$\sim$\SI{3.2}{\volt} \\
Tension nominale & \SI{3.0}{\volt} \\
Capacité nominale & \SI{45}{\milli\ampere\hour} \\
Courant de décharge & \SI{0.12}{\milli\ampere} \\
Dimensions & D : \SI{20.0}{\milli\meter} $\times$ H : \SI{2.0}{\milli\meter} \\
Poids & \SI{2.20}{\gram} \\
Nombre de cycles de vie avant $SoH = 10 \si{\percent}$ & $\approx 1000$ \\
Température de fonctionnement & -20 à \SI{60}{\celsius} \\ 
\bottomrule
\end{tabular}
\end{minipage}\begin{minipage}{0.28\textwidth}
\centering
\includegraphics[width=0.8\textwidth]{Ex24/Fig04}
\end{minipage}
\captionof{table}{Spécifications techniques de l'accumulateur ML2020 de Panasonic{\copyright} [la capacité nominale correspond à la capacité de décharge obtenue lorsque la tension vaut \SI{2.0}{\volt}.]}
\end{center}

D'après le commercial, cette batterie présente une capacité élevée ce qui en fait un dispositif idéal destiné aux clés, télécommandes ou autres objet nécessitant de l'énergie sur le long terme.

L'objectif est d'interpréter quelques expériences ayant conduit à ces spécifications techniques.

\paragraph*{Analyse de la décharge}
Lors de cycles de décharge, un expérimentateur a suivi l'évolution de la tension au sein d'un accumulateur ML2020 neuf dans les conditions expérimentales suivantes : 
\begin{itemize}
\item Charge : un générateur de tension de 3.1 \si{\volt}, une résistance de \SI{160}{\ohm} pendant 60 heures ;
\item Décharge : une résistance de \SI{20}{\kilo\ohm} (équivaut à un courant de décharge de \SI{120}{\micro\ampere}) jusqu'à ce que la tension aux bornes de l'accumulateur vale \SI{2}{\volt}.
\end{itemize}
L'ensemble a été réalisé à 4 températures différentes.

\begin{minipage}{0.48\textwidth}
\includegraphics[width=\textwidth]{Ex24/Fig02}
\end{minipage}\hfill\begin{minipage}{0.48\textwidth}
\captionof{figure}{Courbe de décharge à 4 températures différentes : -20, 0, 20 et \SI{60}{\celsius}.}
\end{minipage}

\Q Représenter schématiquement l'état de charge de la batterie lorsque la capacité est 0, 22 et 44.
\solution{L'état de charge représente le rapport entre la quantité de charges stockées sur la quantité de charges maximale dans la batterie, on peut donc représenter les symboles suivants : 
\begin{center}
\begin{tabular}{cc}
\toprule
Capacité du graphe [\si{\milli\ampere\hour}] & Schéma \\
\midrule
0 & \faBattery[4] \\
22 & \faBattery[2] \\
44 & \faBattery[0] \\
\bottomrule
\end{tabular}\end{center}
}

\Q Comment évolue la quantité maximale de charge stockable dans la batterie lorsque la température augmente. Pourquoi ?
\solution{On lit les valeurs suivantes : 
\begin{center}
\begin{tabular}{ccccc}
\toprule
Température [\si{\celsius}] & -20 & 0 & 20 & 60 \\
\midrule
Capacité [\si{\milli\ampere\hour}] & 40 & 42 & 44 & 44 \\ 
\bottomrule
\end{tabular}
\end{center}
La cinétique chimique est à l'origine de la diminution de la température lors d'un processus rédox de charge ou décharge.
On rappelle que la loi d'Arrhénius s'écrit : 
\begin{align*}
k(T) = A \exp \left( - \frac{E_A}{RT} \right)
\end{align*}
avec $A$ le facteur préexponentiel et $E_A$ l'énergie d'activation (grandeur strictement positive).
Par conséquent, une augmentation en température favorise cinétiquement le processus rédox et permet une charge plus rapide, donc une meilleure capacité.\\
\faWarning ~ Il est important de préciser trois choses supplémentaires aux étudiants : 
\begin{enumerate}
\item lors d'un refroidissement, la résistance interne augmente à cause de la perte de "mobilité" (notion non abordée en cours) des ions, illustrable à l'aide de la formule de Thévenin ;
\item si l'élevation de la température est trop grande, d'autres réactions chimiques indésirables sont aussi favorisés cinétiquement (dégradation de l'électrolyte, réaction rédox des métaux de l'électrode$\ldots$) ce qui peut causer une diminution de la capacité à plus haute température. En effet, une partie du courant de charge servira à alimenter ces processus rédox indésirables ; 
\item lors d'un échauffement, la durée de vie de la batterie diminue (c.-à-d. $SoH$ diminue)
\end{enumerate} 
} 
%La batterie décharge moins de charge à basse température, probablement à cause de la perte de mobilité des ions lithium.}

\Q Lors de la charge et de la décharge, il a été observé que la tension moyenne de charge ($\sim \SI{3.0}{\volt}$) est plus grande que la tension moyenne de décharge. Quel est l'origine de ce processus ? Une explication énergétique est attendue.
%La tension nominale, c'est-à-dire lorsque $SoC=100 ~ \si{\percent}$, correspond approximativement à la valeur qu'on aurait lors d'une charge. Qu'observe-t-on lors de la décharge ? Quel peut-être l'explication ?
\solution{La tension de décharge est plus faible que la tension de charge. En effet, lors de la charge, l'expérimentateur est emmené à devoir fournir plus d'énergie que l'énergie chimique stockable pour permettre au lithium de migrer au sein de l'électrolyte, lui permettre de changer de phase$\ldots$
Lors de la décharge, une partie de l'énergie chimique libérée est perdue pour les mêmes raisons, ce qui résulte en une tension de décharge plus faible que la tension de charge.}

\Q \`{A} partir de quel état de charge $SoC$, la tension de décharge diminue-t-elle ? Pourquoi ? Commenter l'influence de la température.%Commenter l'évolution de la tension de décharge avec la température.
\solution{Tant que la batterie est chargée $SoC > 70 ~ \si{\percent}$, la tension de la pile est relativement constante avec la température. Ce qui traduit un coefficient thermique en \si{\volt\per\kelvin} très faible. On rappelle que le coefficient thermique s'exprime comme :
\begin{align*}
e = - \frac{\Delta_r H}{mF} + \frac{\Delta_r S}{mF} T \iff \left( \frac{\partial e}{\partial T} \right)_{P, \xi} = \frac{\Delta_r S}{mF}
\end{align*}
Pour $SoC < 30 ~ \si{\percent}$, la chute de potentiel est dû à l'augmentation de la résistance interne/surtension de l'accumulateur. Microscopiquement, il devient de plus ne plus difficile pour les ions \ce{Li^+} de quitter l'anode et de s'insérer dans la cathode.
}


\paragraph*{\'{E}volution de l'état de santé de l'accumulateur en fonction du nombre de cycles de charge/décharge}

Dans les mêmes conditions que précédemment, l'accumulateur a été soumis à 10000 cycles de charge/décharge.

\begin{minipage}{0.48\textwidth}
\includegraphics[width=\textwidth]{Ex24/Fig01}
\end{minipage}\hfill\begin{minipage}{0.48\textwidth}
\captionof{figure}{\'{E}volution de l'état de santé d'une batterie après $N$ cycles de charge/décharge.}
\end{minipage}

\Q Sachant qu'un condensateur peut faire des centaines de milliers, voire des millions d cycles, comparer cette technologie à celle de la batterie lithium. 
\solution{
Au bout de 100 cycles, cette batterie lithium a perdu \SI{50}{\percent} de ses capacités de stockage. On est très loin des performances de stockage d'un condensateur.
Cette différence est assez simple à comprendre. Le stockage d'énergie d'un condensateur est d'origine physique (confinement d'un champ électrique entre deux plaques chargées), ce qui n'induit pas de modifications chimiques pourvant être des sources microscopiques d'irréversibilité. Au contraire, les batteries sont basées sur le stockage chimique qui s'accompagne nécessairement d'une multitude de sources d'irréversibilité.
}

\Q Vérifier les spécifications $SoH$ du constructeur.
\solution{On observe qu'après 1000 cycles de charge/décharge, l'état de santé de la pile est de \SI{10}{\percent}. Ceci est en accord avec les informations du constructeur.}

Le constructeur garantit les spécifications fournies pour une tension de charge inférieure à \SI{3.3}{\volt}.

\Q Quel serait l'allure de ce graphique si la charge de l'accumulateur avait lieu à une plus grande tension ? Pourquoi ?
\solution{On doit s'attendre à avoir un déclin de l'état de santé de la batterie bien plus rapide. On peut proposer la courbe suivante : 
\begin{center}
\includegraphics[width=0.6\textwidth]{Ex24/Corr01}
\end{center}
On s'attend à avoir une dégradation des matériaux sous l'effet d'une charge $I$ trop importante. \\
\textbf{Nota Bene : } Vous pouvez commenter l'effet de température sur l'évolution du $SoH$ avec le nombre de cycle.
}

\begin{comment}
\paragraph*{\'{E}volution de la capacité en charge continue à \SI{60}{\celsius}}

La capacité 
\begin{itemize}
\item Charge : un générateur de tension de 3.1 à 3.3 \si{volt}, 
\item Décharge : une résistance de \SI{1.5}{\kilo\ohm} jusqu'à ce que la tension aux bornes de l'accumulateur vale \SI{2}{\volt}.
\end{itemize}


\begin{minipage}{0.48\textwidth}
\includegraphics[width=\textwidth]{Ex24/Fig03}
\end{minipage}\hfill\begin{minipage}{0.48\textwidth}
\captionof{figure}{Charge continue à \SI{60}{\celsius} entre 3.1 et \SI{3.3}{\volt}.}
\end{minipage}
\end{comment}

\end{document}
